\section{معرفی توابع و مثال‌های کاربردی}
\subsection{شروع با نمونهٔ آماده}
تابع
\lr{App\_Init()}
را پس از راه‌اندازی HAL، کلاک و GPIO فقط یک بار و تابع
\lr{App\_Process()}
را مرتب در حلقهٔ اصلی اجرا کنید. تعداد خانه‌ها را با
\lr{APP\_SENSOR\_COUNT}
تعیین کنید؛ برای یک سنسور مقدار یک کافی است. آرایه‌های آدرس، دما و وضعیت با
اندیس یکسان به یک سنسور اشاره می‌کنند. اندیس صفر اولین سنسور است. فقط وضعیت
موفق اجازهٔ استفاده از دما را می‌دهد.

برای میکروی دیگر ابتدا callbackهای پایه و زمان را معرفی کنید. کتابخانه خواندن
کامل مسدودکننده، شروع و خواندن جداگانه و رابط آرایه‌ای با انتظار آزاد دارد.
رابط آرایه‌ای با یک سنسور نیز قابل‌استفاده است.

\subsection{راه‌اندازی باس}
تابع
\lr{OneWire\_Init}
ساختار مقصد باس و تنظیم callbackها را می‌گیرد. اگر توابع الزامی و جفت حفاظت
وقفه معتبر باشند، تنظیمات را کپی و خط را رها می‌کند و مقدار درست می‌دهد. این
تابع سنسور را کشف نمی‌کند. قرارداد callbackها در فصل معماری و نمونهٔ برد در
\lr{Application/Src/app.c}
آمده است.
\begin{LTR}
\begin{lstlisting}
OneWire_Bus bus;
/* callbacks contains your board's GPIO/time functions. */
bool bus_ok = OneWire_Init(&bus, &callbacks);
\end{lstlisting}
\end{LTR}
اگر اینیت ناموفق بود ادامه ندهید. ساختار باس باید تا پایان کار مدیر زنده بماند.
مثال‌های بعدی فرض می‌کنند باس با موفقیت راه‌اندازی شده است.

\subsection{گرفتن و بررسی آدرس}
تابع
\lr{DS18B20\_ReadRom}
باس و آرایهٔ خروجی هشت‌بایتی را می‌گیرد و وضعیت عملیات را می‌دهد. فقط وقتی
دقیقاً یک سنسور متصل است از آن استفاده کنید. تابع
\lr{DS18B20\_IsValidRom}
خانواده و CRC آدرس را بررسی می‌کند ولی حضور فعلی آن سنسور را اثبات نمی‌کند.
\begin{LTR}
\begin{lstlisting}
uint8_t rom[DS18B20_ROM_SIZE];
DS18B20_Status result = DS18B20_ReadRom(&bus, rom);
bool address_ok = (result == DS18B20_STATUS_OK);
\end{lstlisting}
\end{LTR}
برای چند سنسور از کشف خودکار مدیر یا توابع
\lr{OneWire\_SearchFirst}
و
\lr{OneWire\_SearchNext}
با ساختار وضعیت متعلق به برنامه استفاده کنید. هر جست‌وجوی موفق یک آدرس با CRC
معتبر می‌دهد؛ خانوادهٔ سنسور را هم کنترل کنید. ترتیب جست‌وجو شمارهٔ مکان نیست.

\subsection{خواندن کامل یک سنسور}
تابع
\lr{DS18B20\_ReadTemperatureBlocking}
باس، آدرس، بودجهٔ اسمی انتظار میلی‌ثانیه‌ای و اشاره‌گر خروجی را می‌گیرد. تبدیل
را شروع می‌کند، منتظر می‌ماند و دمای تأییدشده را می‌خواند. فقط موفقیت خروجی را
تغییر می‌دهد. برای رزولوشن نامعلوم بودجهٔ ۱۰۰۰ میلی‌ثانیه مناسب است؛ سربار تبادل
به بودجهٔ اسمی اضافه می‌شود.
\begin{LTR}
\begin{lstlisting}
float temperature_c = 0.0f;
if (address_ok) {
    result = DS18B20_ReadTemperatureBlocking(
        &bus, rom, 1000U, &temperature_c);
}
bool temperature_valid = (result == DS18B20_STATUS_OK);
\end{lstlisting}
\end{LTR}
این کد ادامهٔ مثال آدرس است. فقط با اعتبار درست از خروجی استفاده کنید؛ صفر
اولیه یا عدد قبلی اندازه‌گیری موفق نیست. این تابع حدود مدت تبدیل برنامه را
متوقف می‌کند. برای انجام کارهای دیگر حلقه، روش بعدی یا مدیر آرایه را انتخاب کنید.

\subsection{جداکردن شروع تبدیل از خواندن}
تابع
\lr{DS18B20\_StartConversion}
فرمان تبدیل را برای آدرس می‌فرستد؛ آدرس تهی یعنی فرمان عمومی باس. موفقیت یعنی
ارسال فرمان، نه پایان اندازه‌گیری. تابع
\lr{DS18B20\_IsConversionReady}
آماده‌بودن را بررسی می‌کند. تابع
\lr{DS18B20\_ReadTemperature}
نتیجهٔ تبدیل انجام‌شده را با کنترل CRC و بازه می‌خواند؛ تبدیل تازه آغاز نمی‌کند
و فقط در موفقیت خروجی را تغییر می‌دهد.

در STM32 از tick سیستم و در میکروی دیگر از ساعت میلی‌ثانیه‌ای یکنواخت استفاده
کنید. قطعهٔ زیر را یک بار اجرا و متغیرهای وضعیت را بین تکرارهای حلقه حفظ کنید:
\begin{LTR}
\begin{lstlisting}
bool waiting = false;
uint32_t deadline = 0;
temperature_valid = false;
result = DS18B20_StartConversion(&bus, rom);
if (result == DS18B20_STATUS_OK) {
    deadline = HAL_GetTick() + 760U; /* After command. */
    waiting = true;
}
\end{lstlisting}
\end{LTR}
قطعهٔ بعدی را مرتب در حلقهٔ اصلی اجرا کنید:
\begin{LTR}
\begin{lstlisting}
if (waiting &&
    (int32_t)(HAL_GetTick() - deadline) >= 0) {
    waiting = false;
    result = DS18B20_IsConversionReady(&bus)
        ? DS18B20_ReadTemperature(&bus, rom, &temperature_c)
        : DS18B20_STATUS_TIMEOUT;
    temperature_valid = (result == DS18B20_STATUS_OK);
}
/* Other main-loop work continues while waiting. */
\end{lstlisting}
\end{LTR}
میان شروع تبدیل و بررسی آماده‌بودن باس را ریست نکنید و جست‌وجو یا فرمان دیگری
نفرستید. مدیر آرایه این توالی را انجام می‌دهد و معمولاً ساده‌تر است.

\subsection{معرفی بافرها به مدیر آرایه}
تابع
\lr{DS18B20\_ArrayInit}
مدیر، باس، آرایهٔ آدرس قابل‌تغییر، بافر جداگانهٔ کشف، آرایهٔ دما، آرایهٔ وضعیت
و تعداد را دریافت می‌کند. همهٔ بافرها همین ظرفیت را دارند و تا پایان کار زنده
می‌مانند. اشاره‌گر یا تعداد نامعتبر، آدرس تکراری و آدرس غیرصفر نامعتبر نتیجهٔ
نادرست می‌دهند. آدرس صفر یعنی کشف خودکار.
\begin{LTR}
\begin{lstlisting}
#define N 4U /* Use 1U for one sensor. */
static DS18B20_Array sensors;
static DS18B20_Rom addresses[N] = {{0}};
static DS18B20_Rom discovered[N];
static float temperatures[N];
static DS18B20_Status status[N];
bool initialized = DS18B20_ArrayInit(
    &sensors, &bus, addresses, discovered,
    temperatures, status, N);
if (initialized) sensors.sample_interval_ms = 250U;
\end{lstlisting}
\end{LTR}
اینیت را فقط یک بار انجام دهید. فاصلهٔ تنظیم‌شده، استراحت پس از چرخه است و
مدت تبدیل نیست. برای آدرس معلوم آرایه را پیش از اینیت پر کنید؛ ترتیب بایت‌ها
خانواده، شش بایت شناسه و CRC است. مدیر می‌تواند آدرس را عوض کند؛ آرایه را ثابت
تعریف نکنید. بافر ایستا از مصرف زیاد پشته جلوگیری می‌کند.

\subsection{پردازش و دریافت خروجی آرایه}
تابع
\lr{DS18B20\_ArrayProcess}
مدیر و tick فعلی را می‌گیرد. مرتب صدا بزنید. نتیجهٔ درست یعنی چرخهٔ کامل پایان
یافته، حتی اگر بعضی خانه‌ها خطا داشته باشند. نتیجهٔ نادرست یعنی هنوز چرخه کامل
نشده و نشانهٔ خرابی نیست. وضعیت هر خانه را جدا بررسی کنید؛ ترتیب خروجی همان
ترتیب آرایهٔ آدرس است.
\begin{LTR}
\begin{lstlisting}
static float latest_c[N];
static bool valid[N];
if (initialized &&
    DS18B20_ArrayProcess(&sensors, HAL_GetTick())) {
    for (size_t i = 0; i < N; ++i) {
        valid[i] = (status[i] == DS18B20_STATUS_OK);
        if (valid[i]) latest_c[i] = temperatures[i];
    }
}
\end{lstlisting}
\end{LTR}
این تصویر اختیاری، عددهای موفق قبلی را نگه می‌دارد ولی خانهٔ ناموفق را نامعتبر
می‌کند. فقط خانهٔ معتبر را برای نمایش یا کنترل استفاده کنید. خروجی خود کتابخانه
هنگام تبدیل و خطا عدد نامعتبر است.

\subsection{تعویض دستی در خانهٔ مشخص}
تابع
\lr{DS18B20\_ArrayReplace}
مدیر، اندیس مقصد از صفر و آدرس معتبر جدید را می‌گیرد. موفقیت نگاشت را تغییر و
دمای قبلی را نامعتبر می‌کند. شکست می‌تواند از اندیس یا آدرس نامعتبر، آدرس
استفاده‌شده در خانهٔ دیگر یا آزادنبودن مدیر باشد. آدرس جدید را از جست‌وجوی کامل
یا اتصال تک‌سنسوری جداگانه بگیرید؛ روی باس چندسنسوری آدرس را با Read ROM نخوانید.
\begin{LTR}
\begin{lstlisting}
/* replacement_rom was obtained and validated earlier. */
if (sensors.phase == 0U) {
    bool replaced = DS18B20_ArrayReplace(
        &sensors, 1U, replacement_rom); /* Second slot. */
    /* Check replaced, then wait for a valid new sample. */
}
\end{lstlisting}
\end{LTR}
مثال به دو خانه یا بیشتر نیاز دارد و باید از همان task پردازش اجرا شود. تغییر
آدرس دمای تازه نمی‌دهد؛ منتظر نمونهٔ معتبر بعدی باشید. نگاشت بدون ذخیره‌سازی
برنامه پس از ریست باقی نمی‌ماند. تعویض خودکار در فصل بعد توضیح داده شده است.

\subsection{تنظیم و تأیید رزولوشن}
تابع
\lr{DS18B20\_WriteScratchpad}
باس، آدرس، آلارم بالا، آلارم پایین و رزولوشن را می‌گیرد. مثال، رزولوشن ۱۰ بیت و
حدود ۳۰ و ۱۰ درجه را تنظیم می‌کند. تنظیم موقتی است و در EEPROM کپی نمی‌شود.
موفقیت ارسال به‌تنهایی تأیید سنسور آدرس‌دار نیست؛ با توابع خواندن حافظه و
دریافت رزولوشن نتیجه را بررسی کنید:
\begin{LTR}
\begin{lstlisting}
result = DS18B20_WriteScratchpad(
    &bus, rom, 30, 10, DS18B20_RESOLUTION_10_BIT);
uint8_t pad[DS18B20_SCRATCHPAD_SIZE];
if (result == DS18B20_STATUS_OK)
    result = DS18B20_ReadScratchpad(&bus, rom, pad);
bool configured = (result == DS18B20_STATUS_OK) &&
    DS18B20_GetResolution(pad) == DS18B20_RESOLUTION_10_BIT &&
    pad[2] == 30U && pad[3] == 10U;
\end{lstlisting}
\end{LTR}
پیش از نمونه‌گیری یا در زمان آزاد مدیر تنظیم کنید. مدیر همچنان انتظار حداکثر
۱۲ بیت را دارد. تابع
\lr{DS18B20\_GetConversionTimeMs}
برای ۹ تا ۱۲ بیت به‌ترتیب ۹۴، ۱۸۸، ۳۷۵ و ۷۵۰ میلی‌ثانیه می‌دهد؛ در زمان‌بندی
اختصاصی حاشیهٔ مناسب را اضافه کنید.

\subsection{توابع خام و معنی خطاها}
تابع خواندن حافظهٔ موقت نه بایت را اعتبارسنجی می‌کند. فقط پس از موفقیت آن تابع
\lr{DS18B20\_DecodeRaw}
مقدار علامت‌دار را استخراج و بیت‌های تعریف‌نشده را حذف می‌کند؛ تابع
\lr{DS18B20\_RawToCelsius}
آن را به درجه تبدیل می‌کند. برای استفادهٔ معمول تابع خواندن دما را انتخاب
کنید که بازهٔ معتبر را نیز کنترل می‌کند.

موفقیت اجازهٔ استفاده می‌دهد و مشغول یعنی انتظار. نبود دستگاه یعنی حضور در
ریست دیده نشده؛ خطای CRC یعنی دادهٔ غیرقابل‌استفاده؛ پایان مهلت یعنی تبدیل
آماده نشده؛ خروج از بازه یعنی مقدار ناممکن؛ ورودی نامعتبر نیازمند اصلاح آدرس یا
فراخوانی است. سنسور آدرس‌دار قطع‌شده می‌تواند با حضور سنسور دیگر خطای CRC بدهد.
وضعیت را معیار بگیرید، نه صرفاً عدد دما. توابع ریست، بیت و بایت اجزای پروتکل
هستند و معمولاً لایهٔ سنسور آن‌ها را صدا می‌زند؛ پیاده‌سازی آن‌ها در سورس قابل
مطالعه است و در انتهای گزارش تکرار نمی‌شود.
